\chapter{Complexity — Sweep}
%%% generated by the blueprint pipeline from TCSlib.Complexity.TuringMachine.Sweep
%%%%%%%%%%%%%%%%%%%%%%%%%%%%%%%%%%%%%%%%%%%%%%%%%%%%%%%%%%%%%%%%%%%%%%%%%%%%%%%
\section{Overview}

This module provides the generic zipper/transduction layer behind sweep-based
tape simulations (the infrastructure for Arora--Barak's Claim~1.6): a finite
window of a work tape represented as two stacks around a frontier, scanned in
either direction; exact-cost lemmas showing that a local transition rule
realizes a complete forward or backward sweep as a finite-state transduction;
table-valued (indexed) transducers that update one entry per visited cell; and
a bound confining an initialized run's work heads and nonblank cells to the
elapsed-time interval.

%%%%%%%%%%%%%%%%%%%%%%%%%%%%%%%%%%%%%%%%%%%%%%%%%%%%%%%%%%%%%%%%%%%%%%%%%%%%%%%
\section{Declarations}

\begin{definition}[Two-stack tape zipper]
\lean{Turing.FinTM.sweepTape}
\label{Turing.FinTM.sweepTape}
\leanok
Given an integer frontier coordinate $z$ and two finite lists $l$ and $r$ of
possibly-blank symbols over an alphabet $A$, this defines the tape content
$\mathbb{Z} \to A \cup \{\text{blank}\}$ that stores the cells strictly to the
left of $z$ in $l$ (nearest cell first) and the cells at $z$ and to its right
in $r$ (in order); every position beyond the extent of either list reads as
blank.
\end{definition}

\begin{lemma}[Reading the zipper at its frontier]
\lean{Turing.FinTM.sweepTape_read}
\label{Turing.FinTM.sweepTape_read}
\leanok
\uses{Turing.FinTM.sweepTape}
\difficulty{2}
Let $z$ be an integer and let $l$ and $r$ be finite lists of possibly-blank
symbols over an alphabet $A$. The two-stack zipper tape with frontier $z$,
left stack $l$, and right stack $r$ reads, at position $z$, exactly the first
entry of $r$ — or blank when $r$ is empty.
\end{lemma}

\begin{lemma}[Write and move right on the zipper]
\lean{Turing.FinTM.sweepTape_right}
\label{Turing.FinTM.sweepTape_right}
\leanok
\uses{Turing.FinTM.indexedVisit, Turing.FinTM.sweepTape}
\difficulty{4}
Let $z$ be an integer, let $l$ and $r$ be finite lists of possibly-blank
symbols over an alphabet $A$, and let $a$, $b$ be possibly-blank symbols.
Overwriting position $z$ with $b$ in the two-stack zipper tape with frontier
$z$, left stack $l$, and right stack $a \mathbin{::} r$ yields exactly the
zipper tape with frontier $z + 1$, left stack $b \mathbin{::} l$, and right
stack $r$: the written cell transfers from the right stack to the left stack.
\end{lemma}

\begin{definition}[Configuration at a forward sweep frontier]
\lean{Turing.FinTM.sweepCfg}
\label{Turing.FinTM.sweepCfg}
\leanok
\uses{Turing.Cfg, Turing.FinTM.sweepTape}
Given an optional machine state $q$, an input-head position $p$, an integer
frontier $z$, two zipper stacks $l$ and $r$ of possibly-blank symbols, and an
output word, this defines the configuration of a one-work-tape Turing machine
whose work tape is the two-stack zipper with frontier $z$, left stack $l$, and
right stack $r$, whose work head sits at $z$, and whose input-head position
and output are the given ones.
\end{definition}

\begin{definition}[Local sweep action]
\lean{Turing.FinTM.sweepAct}
\label{Turing.FinTM.sweepAct}
\leanok
\uses{Turing.Action}
For a machine state $q$, a possibly-blank symbol $a$, and a direction $d$,
this defines the one-work-tape machine action that overwrites the scanned
work-tape cell with $a$, moves the work head one step in direction $d$, leaves
the input head stationary, appends nothing to the output, and passes to
state $q$.
\end{definition}

\begin{lemma}[One forward sweep step on configurations]
\lean{Turing.FinTM.sweepCfg_right}
\label{Turing.FinTM.sweepCfg_right}
\leanok
\uses{Turing.Action.apply, Turing.FinTM.sweepAct, Turing.FinTM.sweepCfg, Turing.FinTM.sweepTape_right, Turing.moveInputPos_zero}
\difficulty{3}
Consider the forward-sweep configuration of a one-work-tape machine with
(optional) state $q$, input-head position $p$, frontier $z$, left stack $l$,
right stack $a \mathbin{::} r$, and output word $\mathit{out}$. Applying to it
the local sweep action that writes the possibly-blank symbol $b$, moves the
work head right, and enters state $q'$ produces exactly the forward-sweep
configuration with state $q'$, the same input-head position $p$ and output
$\mathit{out}$, frontier $z + 1$, left stack $b \mathbin{::} l$, and right
stack $r$.
\end{lemma}

\begin{definition}[Finite-state transduction of a word]
\lean{Turing.FinTM.sweepFold}
\label{Turing.FinTM.sweepFold}
\leanok
Given a local rule $\mathit{visit} \colon R \times C \to R \times C$ mapping a
control state and a cell to a new state and a rewritten cell, and an initial
state $s$, this defines the left-to-right transduction of a word over $C$: the
word is processed cell by cell, each cell being rewritten by
$\mathit{visit}$ while the control state is threaded through, and the result
records both the final control state and the fully rewritten word.
\end{definition}

\begin{lemma}[A local rule realizes a complete forward sweep]
\lean{Turing.FinTM.sweep_run}
\label{Turing.FinTM.sweep_run}
\leanok
\uses{Turing.Action.apply, Turing.Cfg.workTapeSymbols, Turing.FinTM.sweepAct, Turing.FinTM.sweepCfg, Turing.FinTM.sweepCfg_right, Turing.FinTM.sweepFold, Turing.FinTM.sweepTape_read, Turing.MultiTapeTM, Turing.MultiTapeTM.runFrom, Turing.MultiTapeTM.runFrom_succ_eq_step, Turing.MultiTapeTM.runFrom_zero}
\difficulty{5}
Let $M$ be a one-work-tape Turing machine over alphabet $A$ with states $S$,
and let controller states $R$ and controller cells $C$ be encoded into $S$ and
$A$ by maps $\mathit{state}$ and $\mathit{symbol}$. Suppose a local rule
$\mathit{visit} \colon R \times C \to R \times C$ is implemented by $M$'s
transition table: whenever $M$ is in state $\mathit{state}(s)$ and its work
head reads $\mathit{symbol}(a)$ (for any input symbol), its transition is the
local sweep action that writes the encoding of the rewritten cell, moves
right, and enters the encoding of the new controller state. Then, starting
from the forward-sweep configuration with state $\mathit{state}(s)$, frontier
$z$, left stack $l$, and right stack consisting of the encoded word
$\mathit{as}$ followed by $r$, running $M$ for exactly $|\mathit{as}|$ steps
reaches the forward-sweep configuration whose state encodes the final
controller state of the transduction of $\mathit{as}$, whose frontier is
$z + |\mathit{as}|$, and whose left stack is the encoded rewritten word,
reversed, prepended to $l$, with right stack $r$; the input-head position and
the output word are unchanged.
\end{lemma}

\begin{definition}[Configuration at a backward sweep frontier]
\lean{Turing.FinTM.sweepRevCfg}
\label{Turing.FinTM.sweepRevCfg}
\leanok
\uses{Turing.Cfg, Turing.FinTM.sweepTape}
Given an optional machine state $q$, an input-head position $p$, an integer
frontier $z$, two zipper stacks $l$ and $r$ of possibly-blank symbols, and an
output word, this defines the configuration of a one-work-tape Turing machine
whose work head sits at $z$ and whose work tape is the coordinate reflection
of the two-stack zipper: the cell at coordinate $w$ is the zipper tape with
frontier $-z$, left stack $l$, and right stack $r$, read at $-w$. Thus $r$
stores the cells at the frontier and to its left (in scanning order toward
decreasing coordinates) and $l$ stores the cells strictly to the right.
\end{definition}

\begin{lemma}[One backward sweep step on configurations]
\lean{Turing.FinTM.sweepRevCfg_left}
\label{Turing.FinTM.sweepRevCfg_left}
\leanok
\uses{Turing.Action.apply, Turing.FinTM.sweepAct, Turing.FinTM.sweepRevCfg, Turing.FinTM.sweepTape, Turing.FinTM.sweepTape_right, Turing.moveInputPos_zero}
\difficulty{4}
Consider the backward-sweep configuration of a one-work-tape machine with
(optional) state $q$, input-head position $p$, frontier $z$, left stack $l$,
right stack $a \mathbin{::} r$, and output word $\mathit{out}$. Applying to it
the local sweep action that writes the possibly-blank symbol $b$, moves the
work head left, and enters state $q'$ produces exactly the backward-sweep
configuration with state $q'$, the same input-head position $p$ and output
$\mathit{out}$, frontier $z - 1$, left stack $b \mathbin{::} l$, and right
stack $r$.
\end{lemma}

\begin{lemma}[A local rule realizes a complete backward sweep]
\lean{Turing.FinTM.sweep_run_reverse}
\label{Turing.FinTM.sweep_run_reverse}
\leanok
\uses{Turing.Action.apply, Turing.Cfg.workTapeSymbols, Turing.FinTM.sweepAct, Turing.FinTM.sweepFold, Turing.FinTM.sweepRevCfg, Turing.FinTM.sweepRevCfg_left, Turing.FinTM.sweepTape_read, Turing.MultiTapeTM, Turing.MultiTapeTM.runFrom, Turing.MultiTapeTM.runFrom_succ_eq_step, Turing.MultiTapeTM.runFrom_zero}
\difficulty{5}
Let $M$ be a one-work-tape Turing machine over alphabet $A$ with states $S$,
and let controller states $R$ and controller cells $C$ be encoded into $S$ and
$A$ by maps $\mathit{state}$ and $\mathit{symbol}$. Suppose a local rule
$\mathit{visit} \colon R \times C \to R \times C$ is implemented by $M$'s
transition table: whenever $M$ is in state $\mathit{state}(s)$ and its work
head reads $\mathit{symbol}(a)$ (for any input symbol), its transition is the
local sweep action that writes the encoding of the rewritten cell, moves
left, and enters the encoding of the new controller state. Then, starting
from the backward-sweep configuration with state $\mathit{state}(s)$, frontier
$z$, left stack $l$, and right stack consisting of the encoded word
$\mathit{as}$ followed by $r$, running $M$ for exactly $|\mathit{as}|$ steps
reaches the backward-sweep configuration whose state encodes the final
controller state of the transduction of $\mathit{as}$, whose frontier is
$z - |\mathit{as}|$, and whose left stack is the encoded rewritten word,
reversed, prepended to $l$, with right stack $r$; the input-head position and
the output word are unchanged.
\end{lemma}

\begin{lemma}[Turning the zipper round]
\lean{Turing.FinTM.sweepTape_turn}
\label{Turing.FinTM.sweepTape_turn}
\leanok
\uses{Turing.FinTM.sweepTape}
\difficulty{4}
Let $z$ be an integer and let $l$ and $r$ be finite lists of possibly-blank
symbols over an alphabet $A$. The two-stack zipper tape with frontier $z$,
left stack $l$, and right stack $r$ coincides, as a function on $\mathbb{Z}$,
with the coordinate reflection $w \mapsto$ (zipper tape with frontier
$-(z - 1)$, left stack $r$, and right stack $l$, read at $-w$): reversing the
scanning direction exchanges the two stacks.
\end{lemma}

\begin{lemma}[Transduction of a concatenation]
\lean{Turing.FinTM.sweepFold_append}
\label{Turing.FinTM.sweepFold_append}
\leanok
\uses{Turing.FinTM.sweepFold}
\difficulty{3}
Let $\mathit{visit} \colon R \times C \to R \times C$ be a local rule, $s$ an
initial control state, and $\mathit{as}$, $\mathit{bs}$ two words over $C$.
The finite-state transduction of the concatenation
$\mathit{as} \mathbin{+\!\!+} \mathit{bs}$ starting from $s$ is obtained by
transducing $\mathit{as}$ from $s$, then transducing $\mathit{bs}$ from the
resulting control state: the final state is that of the second transduction,
and the rewritten word is the concatenation of the two rewritten words.
\end{lemma}

\begin{definition}[Table-valued indexed transducer]
\lean{Turing.FinTM.indexedVisit}
\label{Turing.FinTM.indexedVisit}
\leanok
Given a family of local rules $\mathit{visit}_i \colon V \times C \to
V \times C$ indexed by a type $I$ with decidable equality, this defines a
transducer whose control state is a table $s \colon I \to V$ and whose cells
are index-tagged pairs $(i, c)$: on such a pair it applies $\mathit{visit}_i$
to the table entry $s(i)$ and the cell $c$, updates the table at $i$ to the
new value, and outputs $i$ tagged with the rewritten cell. Only the entry
named by the visited cell changes.
\end{definition}

\begin{lemma}[Indexed transduction of a duplicate-free block]
\lean{Turing.FinTM.indexedFold}
\label{Turing.FinTM.indexedFold}
\leanok
\uses{Turing.FinTM.indexedVisit, Turing.FinTM.sweepFold}
\difficulty{5}
Let $\mathit{visit}_i \colon V \times C \to V \times C$ be a family of local
rules indexed by a type $I$ with decidable equality, let
$\mathit{cell} \colon I \to C$ assign a cell to each index, and let
$\mathit{is}$ be a list of indices with no duplicates. Transducing, with the
table-valued indexed transducer started at table $s$, the block that lists
each $i \in \mathit{is}$ tagged with $\mathit{cell}(i)$, yields: the table
that at every index $i$ occurring in $\mathit{is}$ holds the new value of
$\mathit{visit}_i$ applied to the \emph{original} entry $s(i)$ and
$\mathit{cell}(i)$, and elsewhere agrees with $s$; together with the block in
which each tagged cell is replaced by its rewritten value, likewise computed
from the original table. Because no index repeats, each local rule sees the
untouched table entry.
\end{lemma}

\begin{lemma}[Indexed transduction of a complete block]
\lean{Turing.FinTM.indexedFold_block}
\label{Turing.FinTM.indexedFold_block}
\leanok
\uses{Turing.FinTM.indexedFold, Turing.FinTM.indexedVisit, Turing.FinTM.sweepFold}
\difficulty{1}
Let $k \in \mathbb{N}$, let $\mathit{visit}_i \colon V \times C \to V \times
C$ be a family of local rules indexed by $i \in \{0, \dots, k-1\}$, let
$\mathit{cell}$ assign a cell to each index, and let $s$ be an initial table.
Transducing, with the table-valued indexed transducer started at $s$, the
block listing every index $i$ exactly once in increasing order tagged with
$\mathit{cell}(i)$, yields the table $i \mapsto$ (new value of
$\mathit{visit}_i$ on $s(i)$ and $\mathit{cell}(i)$) together with the block
of rewritten tagged cells in the same order.
\end{lemma}

\begin{lemma}[Indexed transduction of a reversed complete block]
\lean{Turing.FinTM.indexedFold_block_reverse}
\label{Turing.FinTM.indexedFold_block_reverse}
\leanok
\uses{Turing.FinTM.indexedFold, Turing.FinTM.indexedVisit, Turing.FinTM.sweepFold}
\difficulty{3}
Let $k \in \mathbb{N}$, let $\mathit{visit}_i \colon V \times C \to V \times
C$ be a family of local rules indexed by $i \in \{0, \dots, k-1\}$, let
$\mathit{cell}$ assign a cell to each index, and let $s$ be an initial table.
Transducing, with the table-valued indexed transducer started at $s$, the
block listing every index exactly once in \emph{decreasing} order tagged with
its cell, yields the table $i \mapsto$ (new value of $\mathit{visit}_i$ on
$s(i)$ and $\mathit{cell}(i)$) together with the correspondingly reversed
block of rewritten tagged cells.
\end{lemma}

\begin{lemma}[An explicit blank at the end of the zipper]
\lean{Turing.FinTM.sweepTape_nil}
\label{Turing.FinTM.sweepTape_nil}
\leanok
\uses{Turing.FinTM.sweepTape}
\difficulty{2}
Let $z$ be an integer and $l$ a finite list of possibly-blank symbols over an
alphabet $A$. The two-stack zipper tape with frontier $z$, left stack $l$, and
empty right stack equals the zipper tape with the same frontier and left stack
whose right stack is the single blank cell.
\end{lemma}

\begin{lemma}[Forward sweep step beyond the stored zone]
\lean{Turing.FinTM.sweepCfg_right_any}
\label{Turing.FinTM.sweepCfg_right_any}
\leanok
\uses{Turing.Action.apply, Turing.FinTM.sweepAct, Turing.FinTM.sweepCfg, Turing.FinTM.sweepCfg_right, Turing.FinTM.sweepTape_nil}
\difficulty{4}
Consider the forward-sweep configuration of a one-work-tape machine with
(optional) state $q$, input-head position $p$, frontier $z$, left stack $l$,
an \emph{arbitrary} (possibly empty) right stack $r$, and output word
$\mathit{out}$. Applying the local sweep action that writes the
possibly-blank symbol $b$, moves the work head right, and enters state $q'$
produces the forward-sweep configuration with state $q'$, the same input-head
position and output, frontier $z + 1$, left stack $b \mathbin{::} l$, and
right stack the tail of $r$; when $r$ is empty, the cell consumed is a blank
beyond the stored zone.
\end{lemma}

\begin{lemma}[Backward sweep step beyond the stored zone]
\lean{Turing.FinTM.sweepRevCfg_left_any}
\label{Turing.FinTM.sweepRevCfg_left_any}
\leanok
\uses{Turing.Action.apply, Turing.FinTM.sweepAct, Turing.FinTM.sweepRevCfg, Turing.FinTM.sweepRevCfg_left, Turing.FinTM.sweepTape_nil}
\difficulty{4}
Consider the backward-sweep configuration of a one-work-tape machine with
(optional) state $q$, input-head position $p$, frontier $z$, left stack $l$,
an \emph{arbitrary} (possibly empty) right stack $r$, and output word
$\mathit{out}$. Applying the local sweep action that writes the
possibly-blank symbol $b$, moves the work head left, and enters state $q'$
produces the backward-sweep configuration with state $q'$, the same
input-head position and output, frontier $z - 1$, left stack
$b \mathbin{::} l$, and right stack the tail of $r$; when $r$ is empty, the
cell consumed is a blank beyond the stored zone.
\end{lemma}

\begin{lemma}[Generating a fixed word at exact cost]
\lean{Turing.FinTM.sweep_generate}
\label{Turing.FinTM.sweep_generate}
\leanok
\uses{Turing.Cfg, Turing.MultiTapeTM, Turing.MultiTapeTM.runFrom, Turing.MultiTapeTM.runFrom_succ_eq_step', Turing.MultiTapeTM.step}
\difficulty{4}
Let $M$ be a one-work-tape Turing machine over alphabet $A$, let $w$ be a word
over $A$ to be written, let $d$ be a fixed integer displacement, and let
$\mathit{cfg}$ be a family of configurations of $M$ indexed by a step counter
$i \in \{0, \dots, |w|\}$, an integer coordinate $z$, and two lists $l$, $r$
of possibly-blank symbols. Assume the local one-step rule: for every
$i < |w|$ and all $z$, $l$, $r$, a single step of $M$ from
$\mathit{cfg}(i, z, l, r)$ yields
$\mathit{cfg}(i+1,\, z + d,\, w_i \mathbin{::} l,\, \mathrm{tail}(r))$, where
$w_i$ is the $i$-th symbol of $w$, pushed as a nonblank cell. Then for every
$n \le |w|$, running $M$ for exactly $n$ steps from $\mathit{cfg}(0, z, l, r)$
yields $\mathit{cfg}$ at counter $n$, coordinate $z + dn$, first list the
first $n$ symbols of $w$, reversed, prepended to $l$, and second list $r$ with
its first $n$ entries removed. The hypothesis is purely local and carries no
premise about the global run.
\end{lemma}

\begin{lemma}[Heads and nonblank cells stay within elapsed time]
\lean{Turing.FinTM.source_bounds}
\label{Turing.FinTM.source_bounds}
\leanok
\uses{Turing.Action.apply, Turing.Cfg.init, Turing.Cfg.inputSymbol, Turing.Cfg.workTapeSymbols, Turing.FinTM, Turing.MultiTapeTM.initCfg, Turing.MultiTapeTM.runFrom, Turing.MultiTapeTM.runFrom_succ_eq_step', Turing.MultiTapeTM.step, Turing.MultiTapeTM.workTapePos_step_le}
\difficulty{5}
Let $M$ be a finite Turing machine over an alphabet $\Gamma$ (a deterministic
multi-tape machine with finitely many states and some number of work tapes),
let $x$ be an input word, and let $t \in \mathbb{N}$. In the configuration
reached by running $M$ for $t$ steps from its initial configuration on $x$:
every work-tape head position lies in the interval $[-t, t]$, and on every
work tape, every cell at a coordinate $z$ with $z < -t$ or $z > t$ is blank.
\end{lemma}
